% ---
% Informações de dados para CAPA e FOLHA DE ROSTO
% ---
\titulo{Título completo do trabalho: subtítulo se for o caso}
\autor{Nome Completo do Autor}
\areaconcentracao{Ciências da Computação.}
%Abaixo outras áreas:
% Educação matemática.
% Matemática pura.
\local{Salvador}
\data{2024} 
\orientador{Prof. Dr. Nome Completo da Pessoa}

%deixe comentada a opção que não se aplica
\dadosorientador{Orientador}
%\dadosorientador{Orientadora}

% ATENÇÃO: Para as titulações utilize:
% Me. para Mestre
% Ma. para Mestra
% Dr. para Doutor
% Dra. para Doutora

% ---
% Informações de dados para FOLHA DE APROVAÇÂO
% ---

\primeiroexaminador{Prof. Dr. Nome Completo da Pessoa}
\dadosprimeiroexaminador{Examinador interno (DCET-I/UNEB)}
\segundoexaminador{Prof. Dra. Nome Completo da Pessoa}
\dadossegundoexaminador{Examinador interno (DCET-I/UNEB)}
\datadefesa{}

%LEIA COM ATENÇÃO AS LINHAS ABAIXO

% Retire o % do início da linha 31 se o seu trabalho deve conter uma seção descrevendo siglas
\newif\ifsiglas
\siglastrue

% Retire o % do início da linha 35 se o seu trabalho deve conter uma seção descrevendo símbolos
\newif\ifsimbolos
%\simbolostrue

% Retire o % do início da linha 39 se o seu trabalho deve conter uma seção de anexos
\newif\ifanexos
%\anexostrue

% Retire o % do início da linha 43 se o seu trabalho deve conter uma seção de apêndices
\newif\ifapendices
%\apendicestrue

% Coloque o % do início da linha 47 se o seu trabalho não deve conter uma página com lista de Figuras
\newif\iffiguras
\figurastrue

% Coloque o % do início da linha 51 se o seu trabalho não deve conter uma página com lista de Tabelas
\newif\iftabelas
\tabelastrue

% Coloque o % do início da linha 53 se o seu trabalho não deve conter uma página com lista de Quadros
\newif\ifquadros
\quadrostrue
